\documentclass{article}
%\usepackage{mathbbol}
\usepackage[utf8]{inputenc}
\usepackage{amsmath}
\usepackage{microtype}
\usepackage{multirow}
\usepackage{graphicx}
\usepackage{float}% If comment this, figure moves to Page 2 %https://tex.stackexchange.com/questions/32598/force-latex-image-to-appear-in-the-section-in-which-its-declared
\usepackage[a4paper,bindingoffset=0.2in,left=0.2in,right=0.2in,top=0.2in,bottom=0.3in,footskip=.25in]{geometry}
\usepackage{hyperref}
\usepackage{adjustbox}
 

\newcommand{\dw}{\Omega} %discrete-time angular frequency
\newcommand{\aw}{\omega} %continuous-time angular frequency
\newcommand{\adoptedbook}{1} %use 1 for Haykin and 2 for Oppenheim

%\newcommand{\dw}{\omega} %discrete-time angular frequency
%\newcommand{\aw}{\Omega} %continuous-time angular frequency
%\newcommand{\adoptedbook}{2} %use 1 for Haykin and 2 for Oppenheim


%From
%http://tex.stackexchange.com/questions/62974/numbering-equations-in-tabular-environment
\newcommand{\tagarray}{%
\mbox{}\refstepcounter{equation}%
$(\theequation)$%
}
%\textrm{~~~}$(\theequation)$%

%\textheight = 592pt
\title{Fourier Analysis: Properties and Pairs}
\author{Prof. Aldebaro Klautau, Universidade Federal do Pará (UFPA), Belém, Brazil}
%\date{1o semestre de 2014}

\begin{document}

\maketitle

\section{Introduction and Adopted Notation}
\begin{itemize}

\item The continuous-time Fourier series and Fourier transform are denoted by the acronyms FS and FT, respectively, while their discrete-time counterparts are DTFS (Discrete-Time Fourier Series) and DTFT (Discrete-Time Fourier Transform).

\item We therefore have four pairs of equations for Fourier analysis: FS, FT, DTFS, and DTFT. All eight equations assume that the analysis is performed over an infinite-duration time interval. In addition to these four pairs, we also have the Discrete Fourier Transform (DFT), which is used for finite-duration signal analysis (i.e., when the signal can be represented by a finite-length array). Fast Fourier Transform (FFT) algorithms are efficient methods for computing the DFT. The DFT is mathematically equivalent to the DTFS, but the DTFS is interpreted as representing a periodic signal of infinite duration, whereas the DFT represents a finite-duration signal.

\item In this text, $\aw$ denotes the continuous-time \emph{angular} frequency, measured in rad/s. Furthermore, $\aw = 2 \pi f$, where $f$ is the continuous-time (``linear'') frequency measured in hertz (Hz).

\item In this text, $\dw$ denotes the discrete-time \emph{angular} frequency, measured in radians (rad).

\item Note that the textbook [2], \emph{Discrete-Time Signal Processing}, 3rd edition, uses $\omega$ to denote discrete-time angular frequency (in radians) and $\Omega$ for continuous-time angular frequency. The textbook [1] by Haykin et al.\ adopts the opposite convention. This table follows the same convention as [\adoptedbook].

\item $x[n]$ and $x(t)$ denote discrete-time and continuous-time signals, respectively, while $X(e^{j\dw})$ and $X(\aw)$ (or $X(j\aw)$) denote their corresponding Fourier transforms (DTFT and FT). The notation $X(f)$ may also be used instead of $X(\aw)$.

\item The symbols $c_k$ and $X[k]$ are used for the continuous-time and discrete-time Fourier series coefficients, respectively.

\item When dealing with periodic discrete-time signals, $N$ samples correspond to the fundamental period, and $\dw_0 = 2\pi/N$ is the fundamental frequency in radians. In continuous time, the period is denoted by $T_0$, and the fundamental frequency can be expressed as either $F_0 = 1/T_0$ Hz or $\aw_0 = 2\pi/T_0$ rad/s.

\item Assuming that $x[n]$ and $x(t)$ are measured in volts, the Fourier series coefficients $c_k$ and $X[k]$ are also measured in volts, whereas the Fourier transform $X(f)$ has units of volts/Hz and can be interpreted as a ``voltage density'' along the frequency axis. An analogy can be drawn between the FS and DTFS and probability mass functions, and between the FT and DTFT and probability density functions. The quantities $X(\aw)/(2\pi)$ and $X(e^{j\dw})/(2\pi)$ may be interpreted as volts per rad/s and volts per rad, respectively, taking into account the factor $1/(2\pi)$ that appears in the corresponding inverse-transform expressions.

\item Fourier transforms are primarily intended to represent non-periodic signals, but they can also represent the spectrum of periodic signals. In this case, each Fourier series coefficient $c_k$ or $X[k]$ appears as the area of the corresponding impulse in the transform domain. For example, a cosine wave with amplitude 16 has Fourier series coefficients $c_1 = c_{-1} = 8$, and its Fourier transform representation is $8\delta(f-F_0)$ and $8\delta(f+F_0)$ when using $X(f)$, or $16\pi\delta(\aw-\aw_0)$ and $16\pi\delta(\aw+\aw_0)$ when using $X(\aw)$.

\item Practice using the linearity property. For example, if the table provides the transform $X(f)$ of a sine wave with unit amplitude, then the transform of a sine wave with amplitude 5 is simply $5X(f)$.

\item There is a duality between periodicity and discreteness. If a signal is periodic in one domain, it will be discrete in the other domain, and vice versa. For example, when a signal is continuous-time and periodic, its spectrum is non-periodic (inherited from the continuous-time nature) and discrete (inherited from the periodicity).

\item The function
\[
\textrm{sinc}(x)=\frac{\sin(\pi x)}{\pi x}
\]
has zeros at $x=\pm 1,\pm 2,\pm 3,\ldots$, and its value at the origin is $\textrm{sinc}(0)=1$.

\end{itemize}

\section{Definitions and Hints for Continuous-Time Fourier Analysis: FS and FT}
\begin{itemize}

\item Switching between $X(f)$ and $X(\aw)$ is often useful. For example, given a plot of $X(f)$, one can obtain the corresponding plot of $X(\aw)$ simply by multiplying the horizontal axis by $2\pi$ (or dividing by $2\pi$ when converting from $X(\aw)$ to $X(f)$), without changing the amplitudes. However, when impulses are involved, a factor of $2\pi$ must be used to scale their areas when converting between $X(f)$ and $X(\aw)$. Due to the impulse-scaling property, an impulse $A\delta(f)$ with area $A$ in $X(f)$ becomes $2\pi A\,\delta(\aw)$ with area $2\pi A$ in $X(\aw)$.

\item The Fourier series coefficients of a periodic time-domain signal can be obtained by sampling, in the frequency domain, the Fourier transform of a single period of the signal. For example, consider the signal $rect_{T_1}(t)$ from Eq.~(\ref{FT2}) as the building block used to generate the periodic pulse train of Eq.~(\ref{FS2}) through convolution in time with the impulse train of Eq.~(\ref{FS6}). In the frequency domain, this convolution corresponds to sampling, since the spectra are multiplied. Consequently,
\[
c_k = \frac{1}{T_0}\,X(f)\bigg|_{f=\frac{k}{T_0}},
\]
which yields Eq.~(\ref{FS2}).

\end{itemize}

\begin{table}[H]
\begin{center}
\caption{Continuous-Time Fourier Transform (FT) Pairs with $X(f)$ Expressed in Hertz}
% Appendix C.4
\begin{tabular}{|c|c|c|c|}

\hline
Eq. & Time Domain & Frequency Domain & Comment \\
\hline
\tagarray\label{FTb}
& $x(t)=\int_{-\infty}^\infty X(f)e^{j 2 \pi f t}df$
& $X(f)=\int_{-\infty}^\infty x(t)e^{-j2 \pi f t}dt$
& frequency expressed in Hz (not rad/s) \\
\hline

\tagarray\label{FT2}
& $\textrm{rect}_{T_1}(t)=\left\{\begin{array}{lr}
1,&|t|\le T_1/2 \\
0,& \textrm{otherwise}
\end{array} \right.$
& $X(f) = T_1 \textrm{sinc}(f T_1)$
& Zeros at $f=\pm 1/T_1, \pm 2/T_1, \cdots$ \\
& & & with ``bandwidth'' $1/T_1$ Hz \\
\hline

\tagarray\label{FT3}
& $\frac{1}{\pi t} \sin(2 \pi F t) = 2 F \textrm{sinc}(2Ft)$
& $X(f)= \left\{\begin{array}{lr}
1,&|f| \le F \\
0,& \textrm{otherwise}
\end{array} \right.$
& pulse $\overset{FT}{\longleftrightarrow}$ sinc. $F$ in Hz \\
\hline

\tagarray\label{FT4}
& $\delta(t)$
& $X(f)=1$
& vertical $\overset{FT}{\longleftrightarrow}$ horizontal \\
\hline

\tagarray\label{FT5}
& $1$
& $X(f)=\delta(f)$
& vertical $\overset{FT}{\longleftrightarrow}$ horizontal \\
\hline

\tagarray\label{FT6}
& $u(t)$
& $X(f) = \frac{1}{j 2 \pi f} +\frac{\delta(f)}{2}$
& $1/(j 2 \pi f)$ is an ``integrator'' \\
\hline

\tagarray\label{FT7}
& $e^{-at}u(t),\hspace{1cm} \Re(a) >0$
& $X(f) = \frac{1}{a+j2\pi f}$
& See Eq.~(\ref{FT7}) \\
\hline

\tagarray\label{FT8}
& $te^{-at}u(t),\hspace{1cm} \Re(a) >0$
& $X(f) = \frac{1}{(a+j2 \pi f)^2}$
& $\int e^{f(x)} dx = f(x)/f'(x)$ \\
\hline

\tagarray\label{FT9}
& $e^{-a|t|},\hspace{1cm}a>0$
& $X(f) = \frac{2a}{a^2+(2 \pi f)^2}$
& Split into the cases $t<0$ and $t \ge 0$ \\
\hline

\tagarray\label{FT10}
& $\frac{1}{\sqrt{2\pi}} e^{-t^2/2}$
& $X(f) = e^{-(2 \pi f)^2/2}$
& Gaussian $\overset{FT}{\longleftrightarrow}$ Gaussian \\
\hline

\end{tabular}
\end{center}
\end{table}

\begin{table}[H]
\begin{center}
\caption{Continuous-Time Fourier Transform (FT) Pairs with $X(\aw)$ Expressed in rad/s (and some also shown using $X(f)$).}
% Appendix C.4
\begin{tabular}{|c|c|c|c|}

\hline
Eq. & Time Domain & Frequency Domain & Comment \\
\hline

\tagarray\label{FT1a}
&
$x(t)=\frac{1}{2{\pi}}\int_{-\infty}^\infty X(j\aw)e^{j\aw t}d\aw$
&
$X(j\aw)=\int_{-\infty}^\infty x(t)e^{-j\aw t}dt$
&
$\aw=2\pi f$, hence $d\aw = 2\pi df$ \\
\tagarray\label{FTbw}
&
or $x(t)=\int_{-\infty}^\infty X(f)e^{j 2 \pi f t}df$
&
$X(f)=\int_{-\infty}^\infty x(t)e^{-j2 \pi f t}dt$
&
and the factor $\frac{1}{2\pi}$ appears in (\ref{FT1a}) \\
\hline

\tagarray\label{FT2w}
&
$\textrm{rect}_{T_1}(t)=\left\{\begin{array}{lr}
1,&|t|\le T_1/2 \\
0,& \textrm{otherwise}
\end{array} \right.$
&
$X(f) = T_1 \textrm{sinc}(f T_1)$
&
Zeros at $f=\pm 1/T_1, \pm 2/T_1, \cdots$ \\
&
&
or $X(j\aw)=\frac{2}{\aw}\sin \left( \frac{\aw T_1}{2} \right)$
&
with ``bandwidth'' $1/T_1$ Hz \\
\hline

\tagarray\label{FT3w}
&
$\frac{1}{\pi t} \sin(Wt)$
&
$X(j\aw)= \left\{\begin{array}{lr}
1,&|\aw| \le W \\
0,& \textrm{otherwise}
\end{array} \right.$
&
pulse $\overset{FT}{\longleftrightarrow}$ sine, NOT sinc. $W$ in rad/s \\
\hline

\tagarray\label{FT4w}
&
$\delta(t)$
&
$X(j\aw)=1$ or $X(f)=1$
&
vertical $\overset{FT}{\longleftrightarrow}$ horizontal \\
\hline

\tagarray\label{FT5w}
&
$1$
&
$X(j\aw)=2\pi\delta(\aw)$ or $X(f)=\delta(f)$
&
vertical $\overset{FT}{\longleftrightarrow}$ horizontal \\
\hline

\tagarray\label{FT6w}
&
$u(t)$
&
$X(j\aw)=\frac{1}{j\aw}+\pi\delta(\aw)$
&
$1/(j\aw)$ is an ``integrator'' \\
\hline

\tagarray\label{FT7w}
&
$e^{-at}u(t),\hspace{1cm} \Re(a)>0$
&
$X(j\aw)=\frac{1}{a+j\aw}$
&
Eq.~(\ref{CTRP3}) with $\gamma=aj$ \\
\hline

\tagarray\label{FT8w}
&
$te^{-at}u(t),\hspace{1cm} \Re(a)>0$
&
$X(j\aw)=\frac{1}{(a+j\aw)^2}$
&
$\int e^{f(x)}dx = f(x)/f'(x)$ \\
\hline

\tagarray\label{FT9w}
&
$e^{-a|t|},\hspace{1cm} a>0$
&
$X(j\aw)=\frac{2a}{a^2+\aw^2}$
&
Split into the cases $t<0$ and $t\ge 0$ \\
\hline

\tagarray\label{FT10w}
&
$\frac{1}{\sqrt{2\pi}}e^{-t^2/2}$
&
$X(j\aw)=e^{-\aw^2/2}$
&
Gaussian $\overset{FT}{\longleftrightarrow}$ Gaussian \\
\hline

\end{tabular}
\end{center}
\end{table}


\begin{table}[H]
\begin{center}
\caption{Continuous-Time Fourier Series (FS) Pairs, where $T_0$ is the fundamental period of $x(t)=x(t+T_0),~\forall t$.}
% Appendix C.2
\begin{tabular}{|c|c|c|c|}

\hline
Eq. & Time Domain & Frequency Domain & Comment \\
\hline

\tagarray\label{FS1}
&
$x(t) = \sum\limits_{k=-\infty}^\infty c_k e^{jk\aw_0t}$
&
$c_k = X[k]= \dfrac{1}{T_0}\displaystyle\int_{\langle T_0 \rangle} x(t)e^{-jk\aw_0t}dt$
&
$\int_{\langle T_0 \rangle}$ denotes integration over one period
\\
\hline

$\vspace{-9pt}$ & \\
\tagarray\label{FS2}
&
$x(t)=  \left\{ {\begin{array}{l}
1, \hspace{6pt}|t| \le T_1/2 \\
0, \hspace{6pt} T_1/2 < |t| \le T_0/2
\end{array}}  \right.$
&
$c_k=\dfrac{\sin \left(k\aw_0 (T_1/2) \right)} {k\pi}
= \frac{T_1}{T_0} \textrm{sinc} \left( \frac{k T_1}{T_0} \right)$
&
Pulse $\overset{FS}{\longleftrightarrow}$ sinc
\\
$\vspace{-10pt}$ & \\
&
$x(t)=x(t+T_0)$
&
\\
$\vspace{-5pt}$ & \\
\hline

\tagarray\label{FS3}
&
$e^{jp\aw_0t}$
&
$c_k = \delta[k-p]$
&
Complex exponential
\\
$\vspace{-9pt}$ & \\
\hline

$\vspace{-9pt}$ & \\
\tagarray\label{FS4}
&
$\cos(p\aw_0t)$
&
$c_k = \frac{1}{2}\delta[k-p] + \frac{1}{2}\delta[k+p]$
&
$\cos(x) = \frac{1}{2}(e^{jx}+e^{-jx})$ and Eq.~(\ref{FS3})
\\
$\vspace{-9pt}$ & \\
\hline

$\vspace{-9pt}$ & \\
\tagarray\label{FS5}
&
$\sin(p\aw_0t)$
&
$c_k = \dfrac{1}{2j}\delta[k-p] - \dfrac{1}{2j}\delta[k+p]$
&
$\sin(x) = \frac{1}{2j}(e^{jx}-e^{-jx})$ and Eq.~(\ref{FS3})
\\
$\vspace{-9pt}$ & \\
\hline

$\vspace{-9pt}$ & \\
\tagarray\label{FS6}
&
$\sum_{p=-\infty}^{\infty} \delta(t-pT)$
&
$c_k = \dfrac{1}{T_0}, \textrm{~~~}\forall k$
&
Impulse train
\\
\hline

\end{tabular}
\end{center}
\end{table}


\begin{table}[H]
\begin{center}
\caption{Fourier Transform Pairs of Periodic Continuous-Time Signals $x(t)=x(t+T_0)$ (FS coefficients become impulses in the FT), where $\aw_0 = 2\pi/T_0$. If $X(f)$ is preferred, the impulse areas in $X(j\aw)$ must be divided by $2\pi$. For example, $X(f)=\delta(f-f_0)$ instead of $X(j\aw)=2\pi\delta(\aw-\aw_0)$, where $\aw_0=2\pi f_0$.}
% Appendix C.5
\begin{tabular}{|c|c|c|}

\hline

Eq. & Periodic Time-Domain Signal & Fourier Transform \\

\hline

\tagarray\label{PCFT1}
&
$x(t) = \sum\limits_{k=-\infty}^{\infty}X[k]e^{jk\aw_0t}$
&
$X(j\aw) = 2\pi \sum\limits_{k=-\infty}^{\infty}X[k]\delta(\aw-k\aw_0)$
or
$X(f) = \sum\limits_{k=-\infty}^{\infty}X[k]\delta(f-kF_0)$
\\

\hline

\tagarray\label{PCFT2}
&
$x(t) = \cos(\aw_0 t)$
&
$X(j\aw) = \pi\delta(\aw-\aw_0)+\pi\delta(\aw+\aw_0)$
\\

\hline

\tagarray\label{PCFT3}
&
$x(t) = \sin(\aw_0 t)$
&
$X(j\aw) = \frac{\pi}{j}\delta(\aw-\aw_0)-\frac{\pi}{j}\delta(\aw+\aw_0)$
\\

\hline

\tagarray\label{PCFT4}
&
$x(t) = e^{j\aw_0 t}$
&
$X(j\aw) = 2\pi\delta(\aw-\aw_0)$
\\

\hline

\tagarray\label{PCFT5}
&
$x(t) = \sum\limits_{n=-\infty}^{\infty}\delta(t-nT_0)$
&
$X(j\aw) = \frac{2\pi}{T_0}\sum\limits_{k=-\infty}^{\infty}\delta\!\left(\aw-k\frac{2\pi}{T_0}\right)$
or
$X(j\aw) = \aw_0 \sum\limits_{k=-\infty}^{\infty}\delta(\aw-k\aw_0)$
\\

\hline

$\vspace{-10pt}$ & \\
\tagarray\label{PCFT6}
&
$x(t)= \left\{
{\begin{array}{l}
1,\hspace{6pt}|t|<T_1/2\\
0,\hspace{6pt}T_1/2<|t|<T_0/2
\end{array}}
\right.$
&
$X(j\aw)=
\sum\limits_{k=-\infty}^{\infty}
\frac{2\sin(k\aw_0(T_1/2))}{k}
\delta(\aw-k\aw_0)$
\\

\hline

\end{tabular}
\end{center}
\end{table}


\begin{table}[H]
\begin{center}
\caption{Properties of Continuous-Time Fourier Representations (FT and FS)}
% Appendix C.7, p. 653 only
\resizebox{1.0\textwidth}{!}{
\begin{tabular}{|c|c|c|c|}
\hline

Eq.& & Fourier Transform & Fourier Series \\

&
&
$x(t) \overset{FT}{\longleftrightarrow} X(j\aw)$
&
$x(t)\overset{FS;\aw_{0}}{\longleftrightarrow}X[k]$
\\

&
&
$y(t) \overset{FT}{\longleftrightarrow}Y(j\aw)$
&
$y(t)\overset{FS;\aw_{0}}{\longleftrightarrow}Y[k]$
\\

&
\textit{Property}
&
&
Period = $T_0$
\\

\hline

\tagarray\label{CTRP1}
&
Linearity
&
$ax(t)+by(t)\overset{FT}{\longleftrightarrow}aX(j\aw)+bY(j\aw)$
&
$ax(t)+by(t)\overset{FS;\aw_{0}}{\longleftrightarrow}aX[k]+bY[k]$
\\

\hline

\tagarray\label{CTRP2}
&
Time shift
&
$x(t-t_{0})\overset{FT}{\longleftrightarrow}e^{-j\aw t_{0}}X(j\aw)$
&
$x(t-t_{0})\overset{FS;\aw_{0}}{\longleftrightarrow}e^{-jk\aw_{0}t_{0}}X[k]$
\\

\hline

\tagarray\label{CTRP3}
&
Frequency shift
&
$e^{j\gamma t}x(t)\overset{FT}{\longleftrightarrow}X(j(\aw-\gamma))$
&
$e^{jk_{0}\aw_{0}t}x(t)\overset{FS;\aw_{0}}{\longleftrightarrow}X[k-k_{0}]$
\\

\hline

\tagarray\label{CTRP4}
&
Time scaling
&
$x(at)\overset{FT}{\longleftrightarrow}\frac{1}{|a|}X\!\left(\frac{j\aw}{a}\right)$
&
$x(at)\overset{FS;a\aw_{0}}{\longleftrightarrow}X[k]$
\\

\hline

\tagarray\label{CTRP5}
&
Time differentiation
&
$\frac{d}{dt}x(t)\overset{FT}{\longleftrightarrow}j\aw X(j\aw)$
&
$\frac{d}{dt}x(t)\overset{FS;\aw_{0}}{\longleftrightarrow}jk\aw_{0}X[k]$
\\

\hline

\tagarray\label{CTRP6}
&
Frequency differentiation
&
$-jtx(t)\overset{FT}{\longleftrightarrow}\frac{d}{d\aw}X(j\aw)$
&
-
\\

\hline

\tagarray\label{CTRP7}
&
Integration / Summation
&
$\int_{-\infty}^{t}x(\tau)d\tau
\overset{FT}{\longleftrightarrow}
\frac{X(j\aw)}{j\aw}+\pi X(j0)\delta(\aw)$
&
-
\\

\hline

\tagarray\label{CTRP8}
&
Convolution
&
$\int_{-\infty}^{\infty}x(\tau)y(t-\tau)d\tau
\overset{FT}{\longleftrightarrow}
X(j\aw)Y(j\aw)$
&
$\int_{\langle T_0\rangle}x(\tau)y(t-\tau)d\tau
\overset{FS;\aw_{0}}{\longleftrightarrow}
T_0X[k]Y[k]$
\\

\hline

\tagarray\label{CTRP9}
&
Modulation
&
$x(t)y(t)
\overset{FT}{\longleftrightarrow}
\frac{1}{2\pi}
\int_{-\infty}^{\infty}
X(jv)Y(j(\aw-v))dv$
&
$x(t)y(t)
\overset{FS;\aw_{0}}{\longleftrightarrow}
\sum_{l=-\infty}^{\infty}X[l]Y[k-l]$
\\

\hline

\tagarray\label{CTRP10}
&
Parseval's theorem
&
$\int_{-\infty}^{\infty}|x(t)|^2dt
=
\frac{1}{2\pi}
\int_{-\infty}^{\infty}|X(j\aw)|^2d\aw$
&
$\frac{1}{T_0}\int_{\langle T_0\rangle}|x(t)|^2dt
=
\sum_{k=-\infty}^{\infty}|X[k]|^2$
\\

\hline

\tagarray\label{CTRP11}
&
\multirow{2}{*}{Duality}
&
\multirow{2}{*}{$X(jt)\overset{FT}{\longleftrightarrow}2\pi x(-\aw)$}
&
$x[n]\overset{DTFT}{\longleftrightarrow}X(e^{j\dw})$
\\
&
&
&
$X(e^{jt})\overset{FS:t}{\longleftrightarrow}x[-k]$
\\

\hline

\tagarray\label{CTRP12}
&
\multirow{4}{*}{Symmetry}
&
$x(t)$ real
$\overset{FT}{\longleftrightarrow}$
$X^*(j\aw)=X(-j\aw)$
&
$x(t)$ real
$\overset{FS:\aw_0}{\longleftrightarrow}$
$X^*[k]=X[-k]$
\\

&
&
$x(t)$ purely imaginary
$\overset{FT}{\longleftrightarrow}$
$X^*(j\aw)=-X(-j\aw)$
&
$x(t)$ purely imaginary
$\overset{FS:\aw_0}{\longleftrightarrow}$
$X^*[k]=-X[-k]$
\\

&
&
$x(t)$ real and even
$\overset{FT}{\longleftrightarrow}$
$\mathrm{Im}\{X(j\aw)\}=0$
&
$x(t)$ real and even
$\overset{FS:\aw_0}{\longleftrightarrow}$
$\mathrm{Im}\{X[k]\}=0$
\\

&
&
$x(t)$ real and odd
$\overset{FT}{\longleftrightarrow}$
$\mathrm{Re}\{X(j\aw)\}=0$
&
$x(t)$ real and odd
$\overset{FS:\aw_0}{\longleftrightarrow}$
$\mathrm{Re}\{X[k]\}=0$
\\

\hline

\end{tabular}
}
\end{center}
\end{table}



\section{Definitions and Hints for Discrete-Time Fourier Analysis: DTFS and DTFT}

\begin{itemize}

\item Recall that there is a duality between periodicity and discreteness. The notation $X(e^{j\dw})$ is intended to emphasize that $\dw$ is an angle and that the spectra of discrete-time signals are always periodic with period $2\pi$ rad (or 360 degrees). The expressions for discrete-time spectra are relatively complicated because they must explicitly represent this periodicity.

\end{itemize}


\begin{table}[H]
\begin{center}
\caption{Discrete-Time Fourier Series (DTFS) Pairs. $N$ is the fundamental period of $x[n]=x[n+N],~\forall n$, $\dw_0=2\pi/N$, and $p \in {\cal Z}$ is an integer used as an auxiliary variable.}
% Appendix C.1
\begin{adjustbox}{max width=\textwidth}
\begin{tabular}{|c|c|c|c|}
\hline
Eq. & Time Domain & Frequency Domain & Comment \\
\hline

\tagarray\label{DTFS1}
&
$x[n] = \sum\limits_{k=\langle N \rangle} X[k]e^{jkn\dw_0}$
&
$X[k]= \dfrac{1}{N}\sum\limits_{n=\langle N \rangle} x[n]e^{-jkn\dw_0}$
&
$\langle N \rangle$ denotes an interval such as $0,1,\ldots,N-1$
\\

\hline

$\vspace{-9pt}$ & \\
\tagarray\label{DTFS2}
&
$x[n] = \left\{
{\begin{array}{l}
1, \hspace{6pt}|n| \le M \\
0, \hspace{6pt} M < |n| \le N/2
\end{array}}
\right.$
&
$X[k]=
\dfrac{
\sin \left(
k \dfrac{\dw_0}{2}(2M+1)
\right)}
{
N \sin \left(
k \dfrac{\dw_0}{2}
\right)
}$
&
Pulse / ``quasi-sinc'' duality
\\

$\vspace{-20pt}$ & \\
&
$x[n]=x[n+N]$
&
\\

$\vspace{-5pt}$ & \\
\hline

$\vspace{-9pt}$ & \\
\tagarray\label{DTFS3}
&
$x[n] = e^{jp\dw_0 n}$
&
$X[k] =
\left\{
\begin{array}{lr}
1, & k=p,p\pm N,p\pm 2N,\ldots \\
0, & \textrm{otherwise}
\end{array}
\right.
$
&
Complex exponential
\\

$\vspace{-9pt}$ & \\
\hline

$\vspace{-9pt}$ & \\
\tagarray\label{DTFS4}
&
$x[n] = \cos(p\dw_0 n)$
&
$X[k] =
\left\{
\begin{array}{lr}
\frac{1}{2}, &
k=\pm p,\pm p\pm N,\pm p\pm 2N,\ldots \\
0, &
\textrm{otherwise}
\end{array}
\right.
$
&
$\cos(x)=\frac{1}{2}(e^{jx}+e^{-jx})$ and Eq.~(\ref{DTFS3})
\\

$\vspace{-9pt}$ & \\
\hline

$\vspace{-9pt}$ & \\
\tagarray\label{DTFS5}
&
$x[n] = \sin(p\dw_0 n)$
&
$X[k] =
\left\{
\begin{array}{lr}
\dfrac{1}{2j}, &
k=p,p\pm N,p\pm 2N,\ldots
\vspace{5pt}
\\
\dfrac{-1}{2j}, &
k=-p,-p\pm N,-p\pm 2N,\ldots
\\
0, &
\textrm{otherwise}
\end{array}
\right.
$
&
$\sin(x)=\frac{1}{2j}(e^{jx}-e^{-jx})$ and Eq.~(\ref{DTFS3})
\\

$\vspace{-9pt}$ & \\
\hline

$\vspace{-9pt}$ & \\
\tagarray\label{DTFS6}
&
$x[n] = 1$
&
$X[k] =
\left\{
\begin{array}{ll}
1, & k=0,\pm N,\pm 2N,\ldots \\
0, & \textrm{otherwise}
\end{array}
\right.
$
&
Vertical/horizontal duality
\\

$\vspace{-9pt}$ & \\
\hline

$\vspace{-9pt}$ & \\
\tagarray\label{DTFS7}
&
$x[n] = \sum_{p=-\infty}^{\infty}\delta[n-pN]$
&
$X[k] = \dfrac{1}{N}, \textrm{~~~}\forall k$
&
Vertical/horizontal duality.
\\
&
&
&
Impulses starting at the origin $n=0$.
\\

\hline

\end{adjustbox}
\end{tabular}
\end{center}
\end{table}


\begin{table}[H]
\begin{center}
\caption{Discrete-Time Fourier Transform (DTFT) Pairs}
% Appendix C.3
\begin{tabular}{|c|c|c|}

\hline
Eq. & Time Domain & Frequency Domain \\
\hline

\tagarray\label{DTFT1}
&
$x[n]=\frac{1}{2{\pi}}\int_{\langle 2{\pi}\rangle}X(e^{j\dw})e^{j\dw n} d\dw$
&
$X(e^{j\dw})=\sum_{n=-\infty}^\infty x[n]e^{-j\dw n}$
\\
\hline

\tagarray\label{DTFT2}
&
$x[n]=\left\{
\begin{array}{lr}
1,& |n| \le M \\
0,& \textrm{otherwise}
\end{array}
\right.$
&
$X(e^{j\dw})=\frac{\sin\!\left[\dw\left(\frac{2M+1}{2}\right)\right]}{\sin(\frac{\dw}{2})}$
\\
\hline

\tagarray\label{DTFT3}
&
$x[n] = \alpha^n u[n],\hspace{1cm}|\alpha|<1$
&
$X(e^{j\dw}) = \frac{1}{1-\alpha e^{-j\dw}}$
\\
\hline

\tagarray\label{DTFT4}
&
$x[n] = \delta[n]$
&
$X(e^{j\dw}) = 1$
\\
\hline

\tagarray\label{DTFT5}
&
$x[n] = u[n]$
&
$X(e^{j\dw}) = \frac{1}{1-e^{-j\dw}} + \pi\sum_{p=-\infty}^\infty \delta(\dw-2\pi p)$
\\
\hline

\tagarray\label{DTFT6}
&
$x[n]=\frac{1}{\pi n}\sin(Wn),\hspace{1cm}0<W \le \pi$
&
$X(e^{j\dw}) =
\left\{
\begin{array}{lr}
1,& |\dw| \le W \\
0,& W<|\dw| \le \pi
\end{array}
\right.$
\\
&
&
$X(e^{j\dw}) \hspace{10 mm} \textrm{has period } 2\pi$
\\
\hline

\tagarray\label{DTFT7}
&
$x[n] = (n+1)\alpha^n u[n]$
&
$X(e^{j\dw}) = \frac{1}{(1-\alpha e^{-j\dw})^2}$
\\
\hline

\end{tabular}
\end{center}
\end{table}


\begin{table}[H]
\begin{center}
\caption{Fourier Transform Pairs of Periodic Discrete-Time Signals (DTFS coefficients become impulses in the DTFT)}
% Appendix C.6
\begin{tabular}{|c|c|c|}
\hline
Eq. & Periodic Time-Domain Signal & Discrete-Time Fourier Transform \\
\hline

\tagarray\label{PDFT1}
&
$x[n] = \sum\limits_{k=\langle N\rangle}X[k]e^{jk\dw_0 n}$
&
$X(e^{j\dw}) = 2\pi\sum\limits_{k=-\infty}^{\infty}X[k]\delta(\dw-k\dw_0)$
\\

\hline

\tagarray\label{PDFT2}
&
$x[n] = \cos(\dw_1 n)$
&
$X(e^{j\dw}) =
\pi\sum\limits_{k=-\infty}^{\infty}
\left[
\delta(\dw-\dw_1-2\pi k)
+
\delta(\dw+\dw_1-2\pi k)
\right]$
\\

\hline

\tagarray\label{PDFT3}
&
$x[n] = \sin(\dw_1 n)$
&
$X(e^{j\dw}) =
\frac{\pi}{j}\sum\limits_{k=-\infty}^{\infty}
\left[
\delta(\dw-\dw_1-2\pi k)
-
\delta(\dw+\dw_1-2\pi k)
\right]$
\\

\hline

\tagarray\label{PDFT4}
&
$x[n] = e^{j\dw_1 n}$
&
$X(e^{j\dw}) =
2\pi\sum\limits_{k=-\infty}^{\infty}
\delta(\dw-\dw_1-2\pi k)$
\\

\hline

\tagarray\label{PDFT5}
&
$x[n] = \sum\limits_{k=-\infty}^{\infty}\delta[n-kN]$
&
$X(e^{j\dw})=
\frac{2\pi}{N}\sum\limits_{k=-\infty}^{\infty}
\delta\!\left(\dw-\frac{2\pi k}{N}\right)$
\\

\hline

\end{tabular}
\end{center}
\end{table}


\begin{table}[H]
\begin{center}
\caption{Properties of Discrete-Time Fourier Representations (DTFT and DTFS)}
% Appendix C.7, p. 654 only
\resizebox{1.0\textwidth}{!}{
\begin{tabular}{|c|c|c|c|}
\hline
Eq. & & Discrete-Time FT & Discrete-Time FS \\
&
&
$x[n] \overset{DTFT}{\longleftrightarrow} X(e^{j\dw})$
&
$x[n] \overset{DTFS:\dw_0}{\longleftrightarrow} X[k]$
\\
&
&
$y[n] \overset{DTFT}{\longleftrightarrow} Y(e^{j\dw})$
&
$y[n] \overset{DTFS:\dw_0}{\longleftrightarrow} Y[k]$
\\
&
\textit{Property}
&
&
Period = $N$
\\

\hline

\tagarray\label{DTRP1}
&
Linearity
&
$ax[n] + by[n] \overset{DTFT}{\longleftrightarrow} aX(e^{j\dw})+bY(e^{j\dw})$
&
$ax[n]+by[n] \overset{DTFS:\dw_0}{\longleftrightarrow} aX[k]+bY[k]$
\\

\hline

\tagarray\label{DTRP2}
&
Time shift
&
$x[n-n_0] \overset{DTFT}{\longleftrightarrow} e^{-j\dw n_0}X(e^{j\dw})$
&
$x[n-n_0] \overset{DTFS:\dw_0}{\longleftrightarrow} e^{-jk\dw_0 n_0}X[k]$
\\

\hline

\tagarray\label{DTRP3}
&
Frequency shift
&
$e^{j\Gamma n}x[n] \overset{DTFT}{\longleftrightarrow} X(e^{j(\dw-\Gamma)})$
&
$e^{jk_0\dw_0 n}x[n] \overset{DTFS:\dw_0}{\longleftrightarrow} X[k-k_0]$
\\

\hline

\tagarray\label{DTRP4}
&
\multirow{2}{*}{Time scaling}
&
$x_z[n] = 0,\ n \neq lp$
&
$x_z[n] = 0,\ n \neq lp$
\\
&
&
$x_z[pn] \overset{DTFT}{\longleftrightarrow} X_z(e^{j\dw lp})$
&
$x_z[pn] \overset{DTFS:p\dw_0}{\longleftrightarrow} pX_z[k]$
\\

\hline

\tagarray\label{DTRP5}
&
Time differentiation
&
-
&
-
\\

\hline

\tagarray\label{DTRP6}
&
Frequency differentiation
&
$-jnx[n] \overset{DTFT}{\longleftrightarrow} \frac{d}{d\dw}X(e^{j\dw})$
&
-
\\

\hline

\tagarray\label{DTRP7}
&
Integration / Summation
&
$\sum\limits_{k=-\infty}^{n} x[k]
\overset{DTFT}{\longleftrightarrow}
\frac{X(e^{j\dw})}{1-e^{-j\dw}}
+
\pi X(e^{j0})
\sum\limits_{k=-\infty}^{\infty}\delta(\dw-2\pi k)$
&
-
\\

\hline

\tagarray\label{DTRP8}
&
Convolution
&
$\sum\limits_{l=-\infty}^{\infty}x[l]y[n-l]
\overset{DTFT}{\longleftrightarrow}
X(e^{j\dw})Y(e^{j\dw})$
&
$\sum\limits_{l=\langle N\rangle}x[l]y[n-l]
\overset{DTFS:\dw_0}{\longleftrightarrow}
NX[k]Y[k]$
\\

\hline

\tagarray\label{DTRP9}
&
Modulation
&
$x[n]y[n]
\overset{DTFT}{\longleftrightarrow}
\frac{1}{2\pi}
\int_{\langle 2\pi\rangle}
X(e^{j\Gamma})Y(e^{j(\dw-\Gamma)})d\Gamma$
&
$x[n]y[n]
\overset{DTFS:\dw_0}{\longleftrightarrow}
\sum\limits_{l=\langle N\rangle}X[l]Y[k-l]$
\\

\hline

\tagarray\label{DTRP10}
&
Parseval's theorem
&
$\sum\limits_{n=-\infty}^{\infty}|x[n]|^2
=
\frac{1}{2\pi}
\int_{\langle 2\pi\rangle}
|X(e^{j\dw})|^2d\dw$
&
$\frac{1}{N}
\sum\limits_{n=\langle N\rangle}|x[n]|^2
=
\sum\limits_{k=\langle N\rangle}|X[k]|^2$
\\

\hline

\tagarray\label{DTRP11}
&
\multirow{2}{*}{Duality}
&
$x[n] \overset{DTFT}{\longleftrightarrow} X(e^{j\dw})$
&
$X[n]\overset{DTFS:\dw_0}{\longleftrightarrow}\frac{1}{N}x[-k]$
\\
&
&
$X(e^{jl}) \overset{FS:t}{\longleftrightarrow} x[-k]$
&
\\

\hline

\tagarray\label{DTRP12}
&
\multirow{4}{*}{Symmetry}
&
$x[n]$ real
$\overset{DTFT}{\longleftrightarrow}$
$X^*(e^{j\dw})=X(e^{-j\dw})$
&
$x[n]$ real
$\overset{DTFS:\dw_0}{\longleftrightarrow}$
$X^*[k]=X[-k]$
\\
&
&
$x[n]$ purely imaginary
$\overset{DTFT}{\longleftrightarrow}$
$X^*(e^{j\dw})=-X(e^{-j\dw})$
&
$x[n]$ purely imaginary
$\overset{DTFS:\dw_0}{\longleftrightarrow}$
$X^*[k]=-X[-k]$
\\
&
&
$x[n]$ real and even
$\overset{DTFT}{\longleftrightarrow}$
$\mathrm{Im}\{X(e^{j\dw})\}=0$
&
$x[n]$ real and even
$\overset{DTFS:\dw_0}{\longleftrightarrow}$
$\mathrm{Im}\{X[k]\}=0$
\\
&
&
$x[n]$ real and odd
$\overset{DTFT}{\longleftrightarrow}$
$\mathrm{Re}\{X(e^{j\dw})\}=0$
&
$x[n]$ real and odd
$\overset{DTFS:\dw_0}{\longleftrightarrow}$
$\mathrm{Re}\{X[k]\}=0$
\\

\hline

\tagarray\label{DTRP13}
&
Conjugation
&
&
$x^*[n]\overset{DTFS:\dw_0}{\longleftrightarrow}X^*[-k]$
\\

\hline

\tagarray\label{DTRP14}
&
Time reversal
&
&
$x[-n]\overset{DTFS:\dw_0}{\longleftrightarrow}X[-k]$
\\

\hline

\end{tabular}
}
\end{center}
\end{table}


\section{References}

\begin{itemize}
    \item $[1]$ \emph{Signals and Systems}, by S. Haykin and B. Van Veen, 2001, Bookman.
    \item $[2]$ \emph{Discrete-Time Signal Processing}, by Oppenheim et al., 3rd edition.
\end{itemize}

\section{Acknowledgments}

Text adapted from textbook [1] by the Signals and Systems class of the first semester of 2014 at UFPA.
It was later improved by the student Carlos Serrato Jr.\ in 2014, and by me and many students since then.


\end{document}

